\documentclass[11pt,letterpaper]{article}
\PassOptionsToPackage{hyphens}{url}
\usepackage[margin=.83in,top=.8in,bottom=.8in,headheight=15pt]{geometry}
\usepackage[T1]{fontenc}
\usepackage{lmodern,amsmath,amssymb,mathtools,bm,booktabs,array,graphicx,tikz}
\usetikzlibrary{arrows.meta,positioning,calc,decorations.pathreplacing,patterns}
\definecolor{ink}{HTML}{17344D}\definecolor{accent}{HTML}{087E8B}\definecolor{pale}{HTML}{EDF5F5}\definecolor{gold}{HTML}{AD6B19}
\providecommand{\dd}{\,\mathrm{d}}\providecommand{\R}{\mathbb{R}}
\providecommand{\E}{\mathbb{E}}\providecommand{\Var}{\operatorname{Var}}
\providecommand{\Cov}{\operatorname{Cov}}\providecommand{\diag}{\operatorname{diag}}
\providecommand{\norm}[1]{\left\lVert#1\right\rVert}
\providecommand{\abs}[1]{\left|#1\right|}
\usepackage{microtype,hyperref,fancyhdr,enumitem,amsthm,sectsty,needspace}
\allsectionsfont{\raggedright}
\hypersetup{colorlinks=true,linkcolor=ink,urlcolor=accent,pdfauthor={Zan Ahmad},pdftitle={Feedback, stability, and hydrostatic pressure / Review}}
\setlength{\parindent}{0pt}\setlength{\parskip}{6pt}\linespread{1.07}
\setlist{itemsep=3pt,topsep=4pt,leftmargin=*}
\emergencystretch=2em\clubpenalty=10000\widowpenalty=10000
\pagestyle{fancy}\fancyhf{}\fancyhead[L]{\small Mathematical Cardiology}\fancyhead[R]{\small Lecture 6}
\fancyfoot[L]{\footnotesize Zan Ahmad}\fancyfoot[R]{\footnotesize\thepage}\renewcommand{\headrulewidth}{.3pt}
\newtheorem{proposition}{Proposition}\newtheorem{theorem}{Theorem}\newtheorem{lemma}{Lemma}
\theoremstyle{definition}\newtheorem{definition}{Definition}\newtheorem{example}{Example}
\numberwithin{equation}{section}
\begin{document}
\section*{Feedback, stability, and hydrostatic pressure --- review sheet}
Use this after attempting the worked examples and practice. Explain each claim in words before calculating.

\section*{Key ideas}
\begin{itemize}
\item Name stimulus, measurement, target, error and action separately.
\item A result depends on what is held fixed.
\item Negative feedback can reduce error, but gain and bounds matter.
\item Hydrostatic pressure is not itself a prediction of circulation flow.
\item Submit the required project idea by this meeting.
\end{itemize}
\section*{Common errors and repairs}
\textbf{Lower resistance always doubles output.} Output changes depend on pump behavior and control assumptions; state what is fixed.

\textbf{More updates always achieve the target.} A rate bound can leave a persistent error.

\textbf{The recurrence models standing because gravity was discussed.} Height and venous storage are absent from the recurrence.

\section*{Retrieve without your notes}
\begin{itemize}
\item Why is the measurement at step n used for rate at step n+\allowbreak{}1?
\item Which units convert stroke volume times rate to L/min?
\item What distinguishes verification from validation?
\end{itemize}
\section*{Readiness check}
\begin{itemize}
\item Separate stimulus, measured response, target, error and feedback action.
\item Calculate cardiac output and pressure in a stated fixed-stroke-volume model.
\item Update a discrete feedback rule by hand and identify actuator saturation.
\item Explain a hydrostatic pressure difference without equating it with net circulation flow.
\item Submit a bounded project idea with a comparison and validation plan.
\end{itemize}
\end{document}
